% === Meta ===
\renewcommand{\title}{Quantum Field Theory II}
\renewcommand{\author}{Gioele Mancino}
\newcommand{\professor}{Prof. Tiziano Peraro}
\newcommand{\faculty}{Department of Physics and Astronomy "A. Righi"}
\newcommand{\course}{Theoretical Physics}
\newcommand{\academicyear}{2025/2026}

% === TODOs management ===
\renewcommand*{\marginfont}{\footnotesize\sffamily}

\newwrite\todofile
\newif\ifshowtodo
\showtodotrue
%\showtodofalse

\newcommand{\TODO}[1]{%
    \ifshowtodo
        \ifx\todofileopen\undefined
            \immediate\openout\todofile=todo.txt
            \global\let\todofileopen\relax
        \fi
        \marginnote{\color{red!70!black}\footnotesize
            \begin{minipage}[t]{2.5cm}
                \raggedright\textbf{TODO:} #1
            \end{minipage}%
        }%
        \immediate\write\todofile{Page \number\value{page}: #1}%
    \fi
}

\AtEndDocument{\ifx\todofileopen\undefined\else\immediate\closeout\todofile\fi}

% === QUESTIONs management ===
\newwrite\questionfile
\newif\ifshowquestions
\showquestionstrue
%\showquestionsfalse

\newcommand{\QUESTION}[1]{%
    \ifshowquestions
        \ifx\questionfileopen\undefined
            \immediate\openout\questionfile=questions.txt
            \global\let\questionfileopen\relax
        \fi
        \marginnote{\color{blue!70!black}\footnotesize
            \begin{minipage}[t]{2.5cm}
                \raggedright\textbf{(?)} #1
            \end{minipage}%
        }%
        \immediate\write\questionfile{Page \number\value{page}: #1}%
    \fi
}

\AtEndDocument{\ifx\questionfileopen\undefined\else\immediate\closeout\questionfile\fi}

% === Custom macros ===
\renewcommand{\dd}{\mathrm{d}}
\newcommand{\ee}{\mathrm{e}}
\newcommand{\ii}{\mathrm{i}}
\newcommand{\RR}{\mathbb{R}}
\newcommand{\CC}{\mathbb{C}}
\newcommand{\NN}{\mathbb{N}}

\renewcommand{\d}[1]{\mathrm{d}#1\,}

\newcommand{\adot}[1]{\dot{\textbf{#1}}}
\newcommand{\ddt}{\frac{\mathrm{d}}{\mathrm{d}t}}
\newcommand{\bs}[1]{\boldsymbol{\mathcal{#1}}}
\newcommand{\avg}[1]{\left\langle #1 \right\rangle}

\newcommand{\dvf}[2]{\frac{\mathrm{d} #1}{\mathrm{d} #2}}
\newcommand{\pdvT}[3]{\left(\frac{\partial #1}{\partial #2}\right)_[#3]}

% === feynman diagram macros ===
\newcommand{\feynLine}{%
    \tikz[baseline=-0.6ex, scale=0.5, thick]{
        \draw (0,0) -- (1.5,0);
        \filldraw (0,0) circle (2pt); \filldraw (1.5,0) circle (2pt);
    }%
}

\newcommand{\feynHoriz}{%
    \tikz[baseline=-0.6ex, scale=0.5, thick]{
        \draw (0,0.4) -- (1,0.4);
        \draw (0,-0.4) -- (1,-0.4);
        \filldraw (0,0.4) circle (2pt); \filldraw (1,0.4) circle (2pt);
        \filldraw (0,-0.4) circle (2pt); \filldraw (1,-0.4) circle (2pt);
    }%
}

\newcommand{\feynVert}{%
    \tikz[baseline=-0.6ex, scale=0.5, thick]{
        \draw (0.2,0.4) -- (0.2,-0.4);
        \draw (0.8,0.4) -- (0.8,-0.4);
        \filldraw (0.2,0.4) circle (2pt); \filldraw (0.2,-0.4) circle (2pt);
        \filldraw (0.8,0.4) circle (2pt); \filldraw (0.8,-0.4) circle (2pt);
    }%
}

\newcommand{\feynCross}{%
    \tikz[baseline=-0.6ex, scale=0.5, thick]{
        \draw (0,0.4) -- (1,-0.4);
        \draw (0,-0.4) -- (1,0.4);
        \filldraw (0,0.4) circle (2pt); \filldraw (1,-0.4) circle (2pt);
        \filldraw (0,-0.4) circle (2pt); \filldraw (1,0.4) circle (2pt);
    }%
}

\newcommand{\feynEight}{%
    \tikz[baseline=-0.6ex, scale=0.5, thick]{
        \draw (0,0) circle (0.5);
        \draw (1,0) circle (0.5);
        \filldraw (0.5,0) circle (2pt);
    }%
}

\newcommand{\feynDumbbell}{%
    \tikz[baseline=-0.6ex, scale=0.65, thick]{
        \draw (0,0) circle (0.4);
        \draw (1.4,0) circle (0.4);
        \draw (0.4,0) -- (1.0,0);
        \filldraw (0.4,0) circle (1.5pt);
        \filldraw (1.0,0) circle (1.5pt);
    }%
}

\newcommand{\feynSunset}{%
    \tikz[baseline=-0.6ex, scale=0.6, thick]{
        \draw (0,0) circle (0.5);
        \draw (-0.5,0) -- (0.5,0);
        \filldraw (-0.5,0) circle (1.5pt);
        \filldraw (0.5,0) circle (1.5pt);
    }%
}

\newcommand{\feynLoop}{%
    \tikz[baseline=-0.6ex, scale=0.6, thick]{
        \draw (0,0) circle (0.5);
        \filldraw (-0.5,0) circle (1.5pt);
        \filldraw (0.5,0) circle (1.5pt);
    }%
}

\newcommand{\feynLineFermion}{%
    \tikz[baseline=-0.6ex, scale=0.5, thick]{
        \draw (0,0) -- (1.5,0);
        \draw[->] (0,0) -- (0.9,0);

        \filldraw (0,0) circle (2pt);
        \filldraw (1.5,0) circle (2pt);
    }%
}

\newcommand{\feynLoopFermion}{%
    \tikz[baseline=-0.6ex, scale=0.6, thick, midarrow/.style={
                postaction={decorate},
                decoration={markings, mark=at position 0.55 with {\arrow{>}}}
            }]{
        \draw[midarrow] (-0.5,0) arc (180:0:0.5);
        \draw[midarrow] (0.5,0) arc (0:-180:0.5);

        \filldraw (-0.5,0) circle (1.5pt);
        \filldraw (0.5,0) circle (1.5pt);
    }%
}

\ExplSyntaxOn
\NewDocumentCommand{\inlinefeynman}{m}
{
    % baseline aggiusta l'altezza del diagramma rispetto al testo
    \tikz[baseline=-0.5ex, x=0.6cm, y=0.6cm, line~width=0.6pt] {

        % l_tmpa_int è la coordinata X di partenza (inizializzata a 0)
        % l_tmpb_int è la coordinata X di arrivo del blocco corrente
        \int_zero:N \l_tmpa_int
        \int_zero:N \l_tmpb_int

        \tl_map_inline:nn { #1 }
        {
            % Ad ogni passo, calcoliamo la coordinata di arrivo (X + 1)
            \int_set:Nn \l_tmpb_int { \l_tmpa_int + 1 }

            \str_case:nn { ##1 }
            {
                { - } {
                        % Propagatore: linea dritta. Nodi a sinistra e a destra.
                        \fill (\int_use:N \l_tmpa_int, 0) circle (1.5pt);
                        \fill (\int_use:N \l_tmpb_int, 0) circle (1.5pt);
                        \draw (\int_use:N \l_tmpa_int, 0) -- (\int_use:N \l_tmpb_int, 0);
                    }
                    { o } {
                        % Anello (Bolla): due archi. Nodi a sinistra e a destra.
                        \fill (\int_use:N \l_tmpa_int, 0) circle (1.5pt);
                        \fill (\int_use:N \l_tmpb_int, 0) circle (1.5pt);
                        \draw (\int_use:N \l_tmpa_int, 0) to[out=90, in=90] (\int_use:N \l_tmpb_int, 0);
                        \draw (\int_use:N \l_tmpa_int, 0) to[out=-90, in=-90] (\int_use:N \l_tmpb_int, 0);
                    }
                    { > } {
                        % Linee entranti (convergono nel nodo a destra)
                        \fill (\int_use:N \l_tmpb_int, 0) circle (1.5pt);
                        \draw (\int_use:N \l_tmpa_int, 0.5) -- (\int_use:N \l_tmpb_int, 0);
                        \draw (\int_use:N \l_tmpa_int, -0.5) -- (\int_use:N \l_tmpb_int, 0);
                    }
                    { < } {
                        % Linee uscenti (divergono dal nodo a sinistra)
                        \fill (\int_use:N \l_tmpa_int, 0) circle (1.5pt);
                        \draw (\int_use:N \l_tmpa_int, 0) -- (\int_use:N \l_tmpb_int, 0.5);
                        \draw (\int_use:N \l_tmpa_int, 0) -- (\int_use:N \l_tmpb_int, -0.5);
                    }
            }
            % Aggiorniamo la coordinata di partenza per il prossimo elemento
            \int_set_eq:NN \l_tmpa_int \l_tmpb_int
        }
    }
}
\ExplSyntaxOff